% Appendix: the conformance kit and campaign. Moved from Section~\ref{sec:conformance}. Sources: the
% library's testing.py (check_backend, check_authorizer), tests/test_testing_kit.py, backends/*,
% authorizers/*, evidence/conformance-campaign-run1.md, evidence/conformance-campaign-run2.md,
% investigations/conformance-run1-triage.md. Planted faults counted in tests/test_testing_kit.py at
% 182bcbc (run 2) and 5d0f414 (run 1).

\section{The conformance kit and campaign in detail}\label{app:conformance}

This appendix gives the kit and the campaign that Section~\ref{sec:conformance} summarizes. The kit,
\texttt{check\_backend} and \texttt{check\_authorizer}, ships with the library so that third-party adapters meet the
same standard, and continuous integration runs it against a real instance of every built-in store except Pinecone,
for which we have no API key, and against the built-in authorizers.

\subsection{Method}\label{sec:conf-method}

\paragraph{Two oracles.}
The reference gives the expected answer $E(\varphi) = \{d \in D : m_d \models \varphi\}$, which the
in-memory backend computes. Because the reference is also code, end-to-end checks add an oracle
written without the library's filter code: $s$ may read $d$ if $s$ is among $d$'s users, or if $d$'s
groups contain the everyone marker or one of $s$'s groups. The retrieval checks use the metadata ACL
authorizer, whose filter and check both evaluate the reference, so the second oracle also tests the
first. Answers are compared as sets, and order is checked as a property, since stores score
differently.

\paragraph{An adversarial generator.}
From a seed, the kit draws corpora of up to 20 documents (by default 25 corpora with four filters
each), each document holding zero to four of seven fields, so missing fields are common. Values are
strings, integers, Booleans, lists of either, nested lists or objects, drawn from confusable pairs:
\texttt{1} and \texttt{"1"}; \texttt{"a"} and \texttt{"A"}; and \'e as one code point (U+00E9) or as
\texttt{e} with a combining accent (U+0301). The ACL documents repeat the pairs in user and group
names (\texttt{bob} and \texttt{Bob}, \texttt{eng} and \texttt{Eng}, a group named \texttt{"1"}). A
pair is equal only for a store that converts types, folds case or normalizes Unicode. The field
\texttt{team.name}, generated beside objects such as \texttt{\{"name": "a"\}} under \texttt{team},
catches stores that read dots as paths, and \texttt{gr\"o\ss e-x} lies outside ASCII. Filters nest
And and Or two levels deep; an In holds zero to three values, so empty In filters occur; about one In
in ten adds 300 values that match nothing, and about one subject in five belongs to 300 extra groups.
A store may refuse a field name it cannot express, returning nothing, but may not answer wrongly.
Every failure reports its seed and example, so that it can be replayed.

\paragraph{What the kit checks.}
For each corpus, \texttt{check\_backend} fills a fresh store and checks (1)~unfiltered and
(2)~filtered searches against $E(\varphi)$: at most $k$ hits, no duplicates, no document outside
$E(\varphi)$, $\min(k,|E(\varphi)|)$ of them, finite scores in non-increasing order, and text and
metadata identical to what was stored, since the re-check decides on that metadata; (3)~that the
asynchronous search returns the same identifiers, order, text, metadata and scores, and refuses the
same filters; (4)~that replacing metadata keeps every text, leaves no old key, creates no document
and changes what filters select as the reference says; and (5)~that retrieval through the layer,
pre-filtered and post-filtered and through both drivers, returns no hit the second oracle denies, no
violation, $\min(k,|R(s)|)$ readable hits, and the same answer from both drivers. Checks (1) and (2)
draw $k$ for each search: it equals the corpus size for one of the two unfiltered searches and for
half of the filtered ones, so that every match must come back and a false denial cannot hide outside
the top $k$, and is otherwise uniform from 1 to the corpus size; check (5) draws $k$ from 1 to 5. The
counts are
required in exact mode. Stores whose index is approximate even on these small corpora (Cassandra, and
Vespa and Redis with HNSW) run in a sound-only mode, in which answers may be short but never wrong; a
store that declares exact search is always checked exactly, because the layer reports its short
answers as complete (Appendix~\ref{sec:sys-complete}). \texttt{check\_authorizer} requires one Boolean
per hit; answers independent of the rest of the batch and its order (all hits, each hit alone, and
the batch reversed); asynchronous answers equal to synchronous ones; and $\mathrm{check}$ agreeing with
$m \models \varphi_s$ on every hit whenever a filter is offered. It empties every cache before each
call, so that a cached answer cannot mask a batch-dependent one. It shows consistency, not a correct
policy: an authorizer that lets everyone read everything passes, and so does a scope narrower than
$R(s)$. OPA, OpenFGA, SpiceDB and Cerbos are additionally compared with a reference answer against
their real servers.

\paragraph{Translations without injection.}
Ten stores (Elasticsearch, OpenSearch, Milvus, MongoDB, Weaviate, Redis, LanceDB, Vespa, Cassandra and
Pinecone) index one field of type-tagged terms $f \triangleright t{:}v$, one for every string or
integer value or list element, where $\triangleright$ is a control character (U+001F) and $t$ a type
letter, as in \texttt{level}$\triangleright$\texttt{i:1}. An In becomes set membership on these
terms, and And and Or become the store's own connectives; equality is exact in type and value,
untagged values (Booleans, floats, nested values) never match, and the metadata itself is stored
unindexed and returned as written. Where terms enter query text (Redis, LanceDB and Vespa), and in
Pinecone, they are hex-encoded, so they contain only digits and the letters a to f. Elsewhere values
are bound parameters (pgvector, Milvus and Cassandra) or data structures of the client API
(Elasticsearch, OpenSearch, Qdrant, Weaviate, Chroma and MongoDB), and FAISS applies the reference in
process. Filter parsing and the translators that write query text are fuzzed with
Atheris~\cite{google2026atheris}: parsing must round-trip or raise, pgvector must bind every field
name and value, and LanceDB and Redis queries must contain only hex-encoded terms.

\paragraph{Testing the tester.}
At the commit that the campaign's second run tested (182bcbc), the kit's own suite plants 18 faulty
backends and six faulty authorizers that the kit must reject (17 and six at the first run's commit,
5d0f414, before the rule for stores that declare exact search),
beside positive controls it must accept: the reference backend, the built-in authorizers, an
approximate backend that misses matches but never errs, and a backend that refuses dotted names. The
backends ignore the filter, alter metadata, ignore $k$, misorder hits, miss matches while checked as
exact or while declaring exact search, diverge on the asynchronous path in six ways, match after case
folding, Unicode normalization or conversion of integers to strings, read dots as paths, or replace
metadata by merging old keys or by creating documents. The authorizers let filter and check
disagree, depend on a hit's position in the batch (alone or behind a cache), or answer differently on
the asynchronous path (alone or behind a cache). Several meta-tests record a gap the kit once had: no
\texttt{1}/\texttt{"1"}, case or Unicode pairs in its data; comparing only the identifiers of filtered
asynchronous searches; skipping the asynchronous search when the synchronous one refused a filter;
and cached answers hiding batch dependence. In this history the kit's blind spots lay in what it
generated and what it compared, never in its oracle.

\subsection{A controlled campaign}\label{sec:conf-campaign}

\paragraph{Setup.}
We ran the kit against current releases of the stores and engines the layer supports, with fixed
versions, seeds and counts, in GitHub-hosted runners. There were 29 targets, of which 28 ran. Thirteen
stores were tested: pgvector 0.8.7 on PostgreSQL 17.11, Redis 8.10.2, Cassandra 5.0.9, Elasticsearch
9.5.5, OpenSearch 3.9.0, Vespa 8.763.13, Qdrant 1.19.2, Milvus 3.0.2, Weaviate 1.39.12, MongoDB Atlas
Local 8.0.30, Chroma 1.5.9, LanceDB 0.40.0 and FAISS 1.15.1, beside the in-memory reference as a
control. The fifteenth store target, Pinecone, needs an API key we did not have and was skipped in
both runs, so the layer's Pinecone adapter is not covered here. The 14
engine targets were OPA 1.21.2, OpenFGA 1.22.0, SpiceDB 1.56.2, Cerbos 0.56.0, Cedar 4.12.2 and Casbin
1.43.0, each with a metadata-ACL policy and a tenant policy, plus Ory Keto 26.2.0 and the in-process
metadata-ACL authorizer as a control. Container image digests are recorded with the results. Each
store target ran 20 seeds of nine data variants and each engine target 20 seeds of seven
(Appendix~\ref{app:kit}). A \emph{run} is one (target, seed, variant) and takes the worst outcome of
its examples: \emph{pass}; a \emph{conformance violation}, the kit's first disagreement with the
reference; a \emph{refusal}, a \texttt{ValueError} by which the layer declines data it cannot
represent; or an \emph{error}, anything else. Engines run in consistency mode, which compares each
engine's filter API with its own check API on every hit.

\paragraph{Run 1.}
The first run (wardowl 2.1.0) completed 4,315 runs: 3,612 passed, 84 were violations, none was
refused and 619 were errors. Both controls passed every run. No target returned a hit outside the
reference, and none altered stored metadata. The violations are of two kinds. \emph{The layer's
Cerbos adapter over-approximated} in 39 of the 40 runs whose ACL fields hold a scalar instead of a
list (the fortieth, one seed of the tenant policy, passed). The adapter's translator maps CEL's
\texttt{in}, \texttt{==} and \texttt{hasIntersection} to $\mathrm{In}$, which also matches a scalar:
for a chunk with \texttt{allowed\_users} equal to the string \texttt{"zo\"e"}, the filter admits the
user, as the reference does, while Cerbos's check evaluates the policy's CEL \texttt{in}, which is
defined on lists and maps only, raises, and denies. The policy is the authority, so the filter
over-approximates; the re-check drops every such hit, and with the re-check off they would be
returned. This is a type gap in the adapter's translator, not a fault of Cerbos, which behaves as it
documents. The project's review had found the gap and the adapter's documentation records it, so the
campaign reproduced it rather than discovered it. \emph{Redis and MongoDB missed matching chunks} in
26 and 19 runs with long and huge
string values, always passing the sound-only retry. Both stores truncate a long term when they index
it and look up the whole term when they search: RediSearch keeps at most 4,096 bytes of a tag value,
which the layer's hex encoding reaches at 2,048 bytes of field and value, and the vector-search
index's filter fields keep 8,191 characters (a reproducer without the layer measured the cut; the
campaign's diagnosis had assumed the 8,181 that MongoDB documents for another field type). A value
equal to the kept prefix of a longer one also matches it, a false admit that the re-check catches;
the reproducers show this, while the campaign's generator holds no such pairs. Neither limit is in
the documentation of the field type the layer uses.

Of the errors, 479 came from OpenFGA and SpiceDB, which refused whole batches with HTTP~400. OpenFGA
accepts at most 100 contextual tuples per request, and one generated subject in five is in more than
300 groups; long identifiers exceeded its size limits. SpiceDB object identifiers admit only a
restricted alphabet, which excludes \texttt{zo\"e} and \texttt{7.0}. The layer failed closed in every
case, but only after a round trip, with an error that did not name the limit. The other 140 errors
came from the stores' own limits, which fail loudly and return nothing: Lucene's 32,766-byte term on
write in OpenSearch (20 runs) and Elasticsearch (16 runs; its other four failed on a query, with
\texttt{x\_content\_parse\_exception}, in the same huge-string variant), Milvus's array-element and JSON
sizes (40 runs), and in 20 runs each, Chroma's rules for metadata keys, which refused the write,
SQLite's expression depth of 1,000 under Chroma, which refused the query, and Qdrant's grammar for
field names. Ory Keto and Casbin ran out of the shared time budget after 72 and 43 of 140 runs, all
passing.

\paragraph{Fixes and run 2.}
Four of these called for a change in the layer's adapters, and we made it: OpenFGA and SpiceDB now refuse
a subject they cannot represent before sending anything, with an error naming the limit, and Redis
and MongoDB store and search a term above their limits as its SHA-256 digest, leaving shorter terms
unchanged. The second run tested the fixed adapters with the harness's budgets split across more
jobs and every finding recorded in full; it changed nothing else, and every outcome matched what the
diagnosis predicted. Of 4,480 runs, 3,822 passed, 39 were violations, 479 were refusals and 140
were errors. Redis and MongoDB passed every long- and huge-string run. The 479 OpenFGA and SpiceDB
runs that had failed with HTTP~400 became refusals by the layer, while the same 21, 21, 16 and 23
runs of the four OpenFGA and SpiceDB targets passed as before. Those refusals are 85.5\% of these
engines' 560 runs. A run takes its worst subject's outcome, so the other subjects of a refused run
still reached the engine, but only 81 runs reached it with every subject: the campaign's evidence on
these two engines is thin, and Section~\ref{sec:engines} tests them on data within their limits.
Keto and Casbin completed all 140 runs each, every one
passing. The 140 store errors were the documented limits above, unchanged. The 39 Cerbos runs
disagreed again, and the harness now attributes each disagreeing hit: in all 280, the pushed-down
filter gave the reference answer and Cerbos's check denied, so the filter over-approximated the
policy as written and never the reverse. No run in either campaign returned a hit outside the
reference or altered metadata.

\subsection{The kit's data variants}\label{app:kit}

Each store target ran 20 seeds of nine variants: the kit's default data (ten corpora per seed) and
eight stress variants (three corpora per seed): corpora of up to 150 documents; filters nested six
deep; disjunctions over 100 and 1,000 values; strings up to 40,000 characters (long and huge);
case and Unicode confusables; numeric strings next to integers; and unusual field names. Each engine
target ran 20 seeds of seven variants: default; case and Unicode; numeric strings; 1,000 groups;
300 chunks per check; long identifiers; and scalar instead of list ACL fields. A run is one (seed,
variant) and takes its worst subject's outcome: pass, a conformance violation, a refusal by the
library (a \texttt{ValueError} for data it cannot represent), or an error. The in-memory reference
store and the in-process metadata-ACL engine are controls.
